\section{Introduction}\label{s:intro}

Streaming video now accounts for the majority of consumer internet traffic~\cite{demandsage2025}, and ABR algorithms play a key role in delivering this content. These algorithms adjust the quality of the playback in response to fluctuating network conditions, aiming to minimize playback interruptions and maintain consistent visual quality to provide a better user experience. In recent years, many ABR algorithms have been proposed, ranging from simple rule-based schemes to complex machine learning models designed to handle challenging network conditions such as variable bandwidth and latency.

Evaluating ABR algorithms under realistic and challenging conditions is essential, particularly in cellular networks. These networks introduce variability in throughput, latency, and packet loss that can reduce video performance. However, evaluating performance directly on real mobile networks is difficult since they lack reproducibility and control. For this reason, emulation tools play a crucial role in evaluating ABR algorithms, by replaying real network traces inside an isolated namespace and preserving key network dynamics. Mahimahi~\cite{mahimahi} is a common emulation tool that has been widely used in previous evaluations. However, recent studies have shown that Mahimahi has some limitations and does not sufficiently capture the dynamic behavior of cellular networks. To address these limitations, CellReplay has been proposed as an improved version of Mahimahi. CellReplay~\cite{cellreplay-nsdi25} offers a higher-fidelity emulation that provides better control over bandwidth and delay, reflecting the dynamics of modern cellular traces with better precision. This thesis uses CellReplay to reassess the performance of six well-known ABR algorithms. We use both emulators to replay the same traces, aiming to observe whether the performance and rankings of ABR algorithms are sensitive to emulator fidelity and how different network traces influence their behavior.

To evaluate ABR algorithms, we use the Puffer live streaming platform. Puffer~\cite{puffer} is an open source system developed at Stanford that supports real video playback with built-in ABR modules and a logging mechanism. It provides a web-based client and a media server that streams multiple TV channels. Puffer’s media servers send data in chunks to the web server to enable video playback, creating a realistic environment to evaluate the behavior of ABR algorithms. Since playback runs over TCP with pre-recorded video content, Puffer offers a practical and reproducible setup for academic research.

The main goal of this thesis is to investigate the following core question; \textit{"To what extent do ABR rankings and behavior change when CellReplay is used, and how are they affected by network trace variability?"}. To address this question, we evaluated ABR algorithms across multiple network traces that differ by generation (4G vs. 5G) and signal strengths. The rest of this thesis is structured as follows. Section~\ref{s:background} introduces the ABR algorithms and evaluation metrics used in our evaluation and explains the emulation tools in more detail. In Section~\ref{s:design}, we provide a high-level overview of our methodology and explain how the different components of the system interact with each together. Sections~\ref{s:evaluation} and~\ref{s:discussion}, present and discuss our results in detail. Finally, in Section~\ref{s:related}, we will compare our work with previous ones.
%%% Local Variables:
%%% mode: latex
%%% TeX-master: "../thesis"
%%% End:
